\documentclass[twocolumn]{article}
% Estilo LaTeX do artigo de lesões periapicais.
% Adaptado do modelo do artigo ROCKFACE-LCCMAT (segmentação de poros).
% Compilar com LuaLaTeX a partir da pasta layout_latex.
\usepackage[
  paperwidth=210mm,
  paperheight=279mm,
  left=14mm,
  right=12mm,
  top=17mm,
  bottom=15mm,
  headheight=9mm,
  headsep=4mm,
  footskip=8mm
]{geometry}

\usepackage{fontspec}
\defaultfontfeatures{Ligatures=TeX,Scale=MatchLowercase}
\setmainfont[
  Path=../fonts/,
  UprightFont=texgyretermes-regular.otf,
  BoldFont=texgyretermes-bold.otf,
  ItalicFont=texgyretermes-italic.otf,
  BoldItalicFont=texgyretermes-bolditalic.otf
]{TeX Gyre Termes}
\setsansfont[
  Path=../fonts/,
  UprightFont=texgyreheros-regular.otf,
  BoldFont=texgyreheros-bold.otf,
  ItalicFont=texgyreheros-italic.otf,
  BoldItalicFont=texgyreheros-bolditalic.otf
]{TeX Gyre Heros}

\usepackage[portuguese]{babel}
\usepackage{microtype}
\usepackage{graphicx}
\graphicspath{{../figuras/}{../../Resultados/figuras/R/}}
\usepackage{booktabs}
\usepackage{array}
\usepackage{tabularx}
\usepackage{ragged2e}
\usepackage{enumitem}
\usepackage{titlesec}
\usepackage{caption}
\usepackage{xcolor}
\usepackage{fancyhdr}
\usepackage{hyperref}
\usepackage{dblfloatfix}
\usepackage{balance}

\definecolor{rockblue}{HTML}{0065A3}
\definecolor{rockgreen}{HTML}{CAD85E}
\definecolor{rockorange}{HTML}{C45A00}
\definecolor{scopegray}{HTML}{F3F3F3}

\hypersetup{
  colorlinks=true,
  linkcolor=rockblue,
  urlcolor=rockblue,
  citecolor=rockblue,
  pdftitle={Detecção de lesões periapicais em radiografias panorâmicas},
  pdfsubject={Random Forest por pixel versus U-Net; figuras geradas em R e Python},
  pdfauthor={William Jota},
  pdfcreator={LuaLaTeX - ML Odonto}
}
\urlstyle{same}

\setlength{\columnsep}{4mm}
\setlength{\parindent}{4.5mm}
\setlength{\parskip}{0pt}
\emergencystretch=1.2em
\clubpenalty=10000
\widowpenalty=10000
\displaywidowpenalty=10000

\AtBeginDocument{\fontsize{8.97pt}{10.71pt}\selectfont}

\titleformat{\section}
  {\sffamily\bfseries\fontsize{10.21pt}{12pt}\selectfont}
  {\thesection.}{0.35em}{}
\titlespacing*{\section}{0pt}{4mm}{0.8mm}
\titleformat{\subsection}
  {\sffamily\bfseries\fontsize{8.72pt}{10.4pt}\selectfont}
  {\thesubsection.}{0.35em}{}
\titlespacing*{\subsection}{0pt}{1.8mm}{0.3mm}
\titleformat{name=\section,numberless}
  {\sffamily\bfseries\fontsize{10.21pt}{12pt}\selectfont}
  {}{0pt}{}

\captionsetup{
  font={small},
  labelfont=bf,
  labelsep=space,
  justification=raggedright,
  singlelinecheck=false,
  skip=1.7mm
}
% babel (português) redefine os nomes no início do documento; reaplicar depois dele
\addto\captionsportuguese{\renewcommand{\figurename}{Fig.}\renewcommand{\tablename}{Tabela}}

% Mais figuras de largura total por página, para evitar páginas só com uma figura
\renewcommand{\dbltopfraction}{0.92}
\renewcommand{\dblfloatpagefraction}{0.72}
\renewcommand{\textfraction}{0.06}
\setcounter{dbltopnumber}{3}

\pagestyle{fancy}
\fancyhf{}
\fancyhead[L]{\sffamily\fontsize{7.97pt}{9pt}\selectfont\color{rockblue}Draft manuscript · not peer reviewed}
\fancyhead[R]{\sffamily\bfseries\fontsize{9.96pt}{11pt}\selectfont Article}
\fancyfoot[L]{\sffamily\fontsize{7.97pt}{9pt}\selectfont\color{rockblue}ML Odonto | draft of 27 October 2026}
\fancyfoot[R]{\sffamily\bfseries\fontsize{7.97pt}{9pt}\selectfont\thepage}
\renewcommand{\headrulewidth}{0.6pt}
\renewcommand{\footrulewidth}{0.6pt}

\fancypagestyle{firstpage}{
  \fancyhf{}
  \fancyhead[L]{\sffamily\bfseries\fontsize{16.94pt}{19pt}\selectfont manuscript draft}
  \fancyhead[R]{\sffamily\bfseries\fontsize{14.94pt}{17pt}\selectfont Article}
  \fancyfoot[L]{\sffamily\fontsize{7.97pt}{9pt}\selectfont\color{rockblue}ML Odonto | draft of 27 October 2026}
  \fancyfoot[R]{\sffamily\bfseries\fontsize{7.97pt}{9pt}\selectfont\thepage}
  \renewcommand{\headrulewidth}{3.2pt}
  \renewcommand{\headrule}{\hbox to\headwidth{\color{rockgreen}\leaders\hrule height \headrulewidth\hfill}}
  \renewcommand{\footrulewidth}{0.6pt}
}

% #1 data, #2 título, #3 autores, #4 resumo, #5 palavras-chave, #6 afiliação
\newcommand{\articlefront}[6]{%
  \begin{minipage}{\textwidth}
    \vspace{4mm}
    \begin{flushright}
      {\sffamily\fontsize{7.97pt}{9pt}\selectfont\color{rockblue}#1}
    \end{flushright}
    \vspace{-1mm}\hrule\vspace{6mm}
    {\sffamily\bfseries\fontsize{25.9pt}{28pt}\selectfont #2\par}
    \vspace{9mm}
    {\sffamily\bfseries\fontsize{9.22pt}{13pt}\selectfont #3\par}
    \vspace{3mm}\hrule\vspace{2mm}
    {\sffamily\fontsize{9.96pt}{14pt}\selectfont\RaggedRight #4\par}
    \vspace{2.5mm}
    {\sffamily\fontsize{7.47pt}{9.4pt}\selectfont #5\par}
    \vspace{3mm}\hrule\vspace{2mm}
    {\sffamily\fontsize{7.47pt}{9.4pt}\selectfont #6\par}
    \vspace{3mm}
  \end{minipage}
}

\newcommand{\widefigure}[4][!t]{%
  \begin{figure*}[#1]
    \centering
    \includegraphics[width=\textwidth,height=#2,keepaspectratio]{#3}
    \caption{#4}
  \end{figure*}
}

\newcommand{\scopenote}[1]{%
  \par\noindent
    \colorbox{scopegray}{%
      \parbox{\dimexpr\columnwidth-2\fboxsep\relax}{%
        \sffamily\color{black!72}\setlength{\parindent}{0pt}%
        \textcolor{rockorange}{\rule{2pt}{2.6em}}\hspace{2mm}#1}}
  \par\vspace{2mm}
}

% Marcador visível de informação que os autores precisam fornecer
\newcommand{\pendente}[1]{\textcolor{rockorange}{\textbf{[#1]}}}

\newcolumntype{Y}{>{\RaggedLeft\arraybackslash}X}
\newcolumntype{L}{>{\RaggedRight\arraybackslash}X}
\newcommand{\tablefont}{\sffamily\fontsize{7.2pt}{8.6pt}\selectfont}
\setlist{nosep,leftmargin=4mm}


\begin{document}
\thispagestyle{firstpage}

\twocolumn[
\articlefront
  {Working draft, 27 October 2026}
  {Detecção de lesões periapicais em radiografias panorâmicas com poucos exemplos locais: Random Forest, U-Net em alta resolução e segmentação em duas etapas guiada por ápices}
  {William Jota\textsuperscript{1}}
  {Lesões periapicais aparecem em radiografias panorâmicas como áreas radiolúcidas pequenas e de contraste baixo, e parte delas não é registrada na leitura de rotina. Avaliamos se um detector automático pode ser construído para o acervo de um serviço que dispunha de apenas 20 panorâmicas com lesões contornadas por especialista (45 lesões) e de 422 panorâmicas sem marcação. Todos os modelos foram treinados com dados públicos: 3.924 panorâmicas de Hanói com lesões delimitadas por caixas e, para localizar ápices radiculares, 634 panorâmicas do DENTEX com os dentes contornados. Comparamos um Random Forest por pixel, uma U-Net com codificador ResNet34 em duas resoluções e perdas, e um método em duas etapas que primeiro localiza os ápices e depois segmenta recortes de alta resolução em torno deles. Cada paciente local foi avaliado por um modelo que não o viu no treinamento. O Random Forest encontrou no máximo 16\% das lesões. A U-Net treinada a 1.400 px encontrou 64\% (29/45), sem falsos positivos. Treinada a 2.000 px com perda de Tversky, que pesa mais a lesão perdida, ela encontrou 69\% no ponto de operação (31/45; IC95\% 54--80\%), também sem falsos positivos, e 76\% (34/45) com no máximo 0,5 falso positivo por imagem. O método em duas etapas pôs 93\% das lesões (42/45) dentro de algum recorte, mas encontrou 64--69\%, com precisão de 76\%. Um classificador por dente, que examina cada ápice com mais contexto e aprende com as caixas de Hanói e os diagnósticos por dente do DENTEX, atingiu AUC de 0,98 nos dados públicos, mas encontrou no máximo 67\% das lesões locais, com 4 falsos positivos por imagem. As abordagens perderam as mesmas lesões: 13 das 45, pequenas e sobretudo maxilares, não foram encontradas por nenhum modelo. O gargalo, portanto, não é localizar a região nem a resolução, e sim reconhecer esse tipo de lesão, o que depende de exemplos de treinamento. Aplicada às 422 panorâmicas, a U-Net de 2.000 px assinalou 807 lesões em 407 imagens. Dados públicos com rótulos grosseiros bastam para detectar lesões periapicais evidentes em outro serviço; as lesões pequenas e maxilares exigem marcações locais.}
  {\textbf{Palavras-chave:} lesão periapical; periodontite apical; radiografia panorâmica; segmentação; U-Net; perda de Tversky; localização de ápices; Random Forest.}
  {\textsuperscript{1} Laboratório de Estudos Eletromagnéticos (LEMAG), Universidade de Brasília, Brasília, Brasil.}
]

\section{Introdução}

A periodontite apical é a resposta inflamatória dos tecidos ao redor do ápice radicular à infecção do sistema de canais, e sua manifestação radiográfica típica é uma área radiolúcida junto à ponta da raiz. Em escala populacional, a radiografia panorâmica é o exame mais disponível para esse rastreamento, mas sobreposição de estruturas, distorção e baixo contraste reduzem sua sensibilidade em relação à tomografia computadorizada de feixe cônico [1]. Mesmo especialistas deixam passar parte das lesões: em uma comparação com 24 cirurgiões bucomaxilofaciais, a taxa média de verdadeiros positivos na leitura de panorâmicas foi de 0,51 [2].

Redes neurais convolucionais vêm sendo usadas para apoiar essa leitura. Ekert e colaboradores mostraram detecção de lesões apicais em panorâmicas com uma rede convolucional [3], e Endres e colaboradores relataram desempenho comparável ao de parte dos especialistas com 2.902 panorâmicas rotuladas [2]. Esses estudos dependem de milhares de imagens marcadas no próprio serviço. Clínicas e grupos de pesquisa menores raramente dispõem desse volume, e a marcação densa de lesões por especialistas é cara.

Quem tem poucos exemplos locais tem três vias. A primeira é a classificação por pixel com atributos predefinidos e Random Forest [4], popularizada por ferramentas interativas como ilastik e Trainable Weka Segmentation [5,6], que aprendem com poucos exemplos e são interpretáveis. A segunda é treinar uma rede convolucional em um conjunto público de outro serviço e aplicá-la ao acervo local [7,8]. A terceira explora a anatomia: como a lesão periapical está, por definição, junto a um ápice radicular, um primeiro modelo pode localizar os ápices e um segundo examinar, em resolução maior, apenas a vizinhança de cada um. Essa decomposição reduz a área de busca e deixa a lesão maior em relação ao recorte, o que poderia favorecer lesões pequenas.

Este estudo compara as três vias no acervo de um serviço que tem 20 panorâmicas contornadas e 422 sem marcação. Todos os modelos foram treinados apenas com dados públicos: 3.924 panorâmicas de Hanói com lesões delimitadas por caixas [7] e, para os ápices, 634 panorâmicas do DENTEX com os dentes contornados e numerados [8]. Avaliamos um Random Forest por pixel; uma U-Net [9] em duas configurações, a segunda em resolução maior e com perda de Tversky [18]; o ajuste fino com as imagens locais; o método em duas etapas; e um classificador que decide, dente a dente, se o ápice tem lesão. Todos seguem o mesmo protocolo de avaliação por lesão. Em seguida, caracterizamos que lesões nenhum modelo encontra e aplicamos o modelo escolhido ao acervo sem marcação. A contribuição é uma comparação justa entre arquiteturas no cenário de poucos dados e a identificação do gargalo que limita todas elas.

\widefigure{108mm}{fig3_diagrama.pdf}{\textbf{Fluxo de trabalho.} Dados, pré-processamento, modelos, validação e aplicação. Os modelos da área azul foram treinados no Google Colab apenas com os conjuntos públicos de Hanói e DENTEX; o Random Forest e o ajuste fino foram treinados localmente. As radiografias de pacientes foram processadas apenas no computador local. Na aplicação, o acervo sem marcação passou pelo mesmo pré-processamento e pela U-Net v2.}

\section{Materiais e métodos}

\subsection{Radiografias locais e marcações}

O acervo local contém 440 radiografias panorâmicas digitais em tons de cinza (438 em JPEG e 2 em TIFF), com largura entre cerca de 2.400 e 2.900 pixels. Vinte pacientes tiveram as lesões periapicais contornadas por \pendente{número e formação dos examinadores} em um documento PDF, uma página por paciente; dezoito dessas radiografias também estão no acervo e duas existem apenas no PDF. As 422 radiografias restantes, sem marcação, formaram o conjunto de aplicação.

Os contornos estavam gravados no PDF como desenhos vetoriais sobre a radiografia. Reconstruímos cada radiografia unindo as faixas de imagem da página e rasterizamos os contornos nas mesmas coordenadas, o que dispensou registro entre imagem e marcação. Muitos traços eram arcos abertos em forma de ``$\cap$'', contornando a lesão e excluindo a raiz. Fechamos cada arco ligando suas extremidades livres e preenchemos o interior; quando não havia interior, usamos o fecho convexo do traço. Componentes com menos de 12 pixels de lado foram descartados como resquícios do traço. O resultado foram 45 lesões em 20 pacientes (Fig.~2a), 20 delas na metade superior da imagem, tratada aqui como maxila, e 25 na inferior, tratada como mandíbula.

\subsection{Conjuntos públicos de treinamento}

\textbf{Hanói.} Usamos o conjunto de panorâmicas com lesões periapicais da Universidade de Medicina de Hanói [7], distribuído sob licença CC BY 4.0 e composto por 3.924 radiografias. Cada lesão está delimitada por uma caixa justa ao redor da área radiolúcida, com largura mediana de 60 pixels, e recebe uma classe 3, 4 ou 5, que pode corresponder ao índice periapical [10] e não foi usada. Como a caixa contém lesão e osso adjacente, cada caixa foi convertida em uma elipse inscrita com semieixos a 85\% da meia-caixa. Os pixels entre essa elipse e uma elipse a 150\% foram excluídos da perda, porque podem ser lesão ou fundo (Fig.~2b).

\textbf{DENTEX.} O conjunto DENTEX [8] (licença CC BY-NC-SA 4.0) traz 634 panorâmicas em que cada dente está contornado por um polígono e numerado por quadrante e posição (18.095 dentes). O DENTEX também traz 705 panorâmicas com diagnósticos por dente (cárie, cárie profunda, lesão periapical e dente incluso), com 158 lesões periapicais em 116 imagens. Nesses achados o polígono contorna o dente inteiro, e não a lesão; por isso eles não foram usados como máscara de lesão, apenas como rótulo por dente no classificador descrito adiante. Usamos os polígonos de numeração para obter ápices. Em cada dente, consideramos a faixa de 12\% da altura do polígono mais próxima da raiz (para cima nos quadrantes superiores, para baixo nos inferiores). O ápice foi o ponto extremo dessa faixa. Nos molares, quando a faixa tinha dois extremos separados por mais de 25\% da largura do dente, registramos dois ápices (Fig.~2c).

\begin{figure*}[!t]
\centering
\includegraphics[width=\textwidth]{fig2_rotulos.png}
\caption{\textbf{Rótulos usados no treinamento e na avaliação.} \textbf{a}, marcação de um especialista no PDF (traço) e a máscara derivada após o fechamento do arco e o preenchimento; a raiz contornada pelo traço fica fora da máscara. \textbf{b}, panorâmica do conjunto de Hanói com a caixa convertida em elipse (lesão) e o anel ao redor excluído do treinamento. \textbf{c}, panorâmica do DENTEX com os polígonos dos dentes (verde) e os ápices derivados da ponta de cada raiz (ciano), usados para treinar o localizador de ápices.}
\end{figure*}

\subsection{Pré-processamento}

As radiografias foram redimensionadas, preservando a proporção, e tiveram o contraste equalizado por CLAHE (limite 2,0; grade de $8 \times 8$) [11], que aproxima o contraste de aparelhos diferentes. A escala de trabalho tem 1.400 pixels de largura; a U-Net v2 trabalhou a 2.000 px e os recortes da segunda etapa foram tirados a 2.800 px. Nessa escala de trabalho, a lesão marcada típica tem diâmetro equivalente mediano de 32 pixels (intervalo interquartil 20--45). Medidas de tamanho são dadas em pixels dessa escala, porque as imagens não continham calibração espacial. Todos os mapas de probabilidade foram levados à escala de trabalho antes da avaliação.

\subsection{Random Forest por pixel}

Cada pixel foi descrito por 38 atributos: intensidade; suavizações gaussianas ($\sigma = 1, 2, 4, 8, 16$ px); diferenças de gaussianas (1--4, 2--8, 4--16, 8--32 e 32--64 px); magnitude do gradiente ($\sigma = 1, 2, 4, 8$); autovalores da Hessiana ($\sigma = 2, 4, 8$); desvio-padrão local ($\sigma = 4, 8$); contexto, dado pela diferença entre a intensidade suavizada ($\sigma = 4, 8$) a 20 e 40 px acima e abaixo e a 25 px para cada lado e a do próprio pixel; e posição horizontal, vertical e distância horizontal ao centro.

Sorteamos 200 pixels de lesão e 1.000 de fundo em cada uma de 800 panorâmicas de Hanói (959.989 pixels; 16,7\% de lesão), metade do fundo na faixa central da imagem, onde estão raízes e espaços interdentais. O Random Forest tinha 100 árvores, profundidade máxima 20, mínimo de 10 amostras por folha, $\sqrt{d}$ atributos por divisão e 200.000 pixels por árvore. Aplicamos esse modelo a outras 300 panorâmicas e acrescentamos até 600 pixels de fundo por imagem onde a probabilidade de lesão era $\geq 0{,}3$ (180.000 negativos difíceis); o modelo retreinado foi chamado RF B. A variante RF C somou ao RF B, em validação cruzada, 2.000 pixels de lesão e 8.000 de fundo de cada imagem local de treino.

\subsection{U-Net v1 e ajuste fino}

A U-Net [9] usou como codificador uma ResNet34 [12] pré-treinada no ImageNet, com a imagem em tons de cinza replicada nos três canais, e um decodificador com cinco estágios de duas convoluções $3 \times 3$, normalização em lote e conexões de atalho. O treinamento usou 3.724 panorâmicas de Hanói, reservando 200 para validação interna. A cada passo, um lote de 8 recortes de $512 \times 768$ pixels foi sorteado, metade centrada em uma lesão, com espelhamento horizontal e variações de gama, ganho e brilho. A perda foi a soma da entropia cruzada binária com a perda Dice, ambas calculadas fora do anel excluído. Usamos AdamW [13] com taxa de aprendizado $3 \times 10^{-4}$, decaimento $10^{-4}$, aquecimento de 200 passos e decaimento cossenoidal ao longo de 6.000 passos, em precisão mista, em uma GPU NVIDIA T4 do Google Colab (cerca de 0,45 s por passo). Salvamos o modelo de maior Dice na validação interna (passo 5.500). Na previsão, a probabilidade foi a média entre a imagem e sua versão espelhada.

Para testar se as 20 imagens locais acrescentavam informação, ajustamos essa U-Net com os contornos dos especialistas, em validação cruzada. Para caber em CPU e limitar a memorização, congelamos o codificador e ajustamos apenas o decodificador, por 150 passos com lotes de 2 recortes de $384 \times 576$ pixels, 70\% deles centrados em lesão. Uma faixa de 4 pixels em cada lado do contorno marcado foi excluída da perda, porque o traço manual não é exato. O mesmo procedimento foi aplicado depois à U-Net v2, também em validação cruzada.

\subsection{U-Net v2: resolução maior e perda de Tversky}

A análise da U-Net v1 mostrou que as lesões perdidas eram pequenas (ver Resultados). Treinamos então uma segunda U-Net, com a mesma arquitetura, sobre as panorâmicas de Hanói a 2.000 px de largura, o que aumenta cada lesão em 43\% de diâmetro. A perda Dice foi substituída pela perda de Tversky [18], que pondera separadamente os pixels marcados a mais ($\alpha = 0{,}3$) e os de lesão não marcados ($\beta = 0{,}7$) e, assim, penaliza mais a lesão perdida. Os demais parâmetros foram os da v1, com 12.000 passos; o melhor Dice na validação interna (0,690) ocorreu no passo 10.000. Na previsão, a imagem foi processada a 2.000 px e o mapa de probabilidade reduzido à escala de trabalho.

\subsection{Método em duas etapas guiado por ápices}

\textbf{Etapa 1: localizador de ápices.} Uma U-Net com a mesma arquitetura foi treinada a prever um mapa de calor com uma gaussiana ($\sigma = 6$ px) em cada ápice derivado do DENTEX, na escala de trabalho. Usamos 584 panorâmicas para treinamento e 50 para validação, 3.000 passos com lotes de 8, e entropia cruzada binária com peso 21 nos pixels de ápice para compensar sua raridade. Os ápices previstos foram os máximos locais do mapa acima de 0,3, separados por ao menos 10 px, até 60 por imagem. Nas 50 imagens de validação, 52\% dos ápices anotados tiveram um ápice previsto a até 20 px, e 57\% dos ápices previstos tinham um anotado a essa distância (F1 0,545 no passo 2.500, o modelo salvo). O localizador detectou, em mediana, 30 ápices por panorâmica de Hanói (116.668 em 3.924 imagens) e 30 por panorâmica local (12--39).

\textbf{Etapa 2: segmentador de recortes.} Em cada panorâmica de Hanói, recortamos janelas de $256 \times 256$ pixels a 2.800 px de largura, centradas nos ápices detectados, o que equivale a $128 \times 128$ pixels na escala de trabalho. Guardamos até 8 recortes com lesão por imagem, nos ápices mais próximos das lesões, e 4 sem lesão. Ao todo foram 14.734 recortes com lesão e 14.882 sem lesão para treinamento, e 765 e 800 para validação, de imagens separadas. Uma U-Net ResNet34 foi treinada nesses recortes por 6.000 passos, com lotes de 32, deslocamento aleatório de até 16 px e perda BCE + Tversky; o melhor Dice nos recortes de validação foi 0,771 (passo 6.000). Na previsão, cada panorâmica local passou pelo localizador, o segmentador foi aplicado a todos os recortes (com espelhamento), e os recortes sobrepostos foram combinados pelo máximo. Fora dos recortes, a probabilidade foi zero.

\subsection{Classificador por dente}

Para testar se mais contexto em torno de cada ápice e rótulos de outro serviço ajudariam a reconhecer as lesões, treinamos um classificador que responde, para cada ápice detectado, se há lesão periapical. Cada ápice originou um recorte de $384 \times 384$ pixels a 2.800 px de largura ($192 \times 192$ na escala de trabalho, o dobro da área dos recortes da segunda etapa); recortes da metade inferior da imagem foram espelhados na vertical, para que a raiz apontasse sempre para cima e maxila e mandíbula compartilhassem o aprendizado. Em Hanói, um ápice a até 40 px de uma caixa de lesão foi rotulado como positivo e um ápice a mais de 90 px de todas, como negativo; os intermediários foram ignorados. No DENTEX com diagnósticos, os ápices do dente com lesão periapical foram positivos e os demais ápices da imagem, negativos. O conjunto teve recortes de 3.924 panorâmicas de Hanói e 705 do DENTEX, com até 6 e 12 negativos por imagem, respectivamente. Uma ResNet34 [12] pré-treinada no ImageNet foi treinada por 5.000 passos (lotes de 48, entropia cruzada binária, AdamW), com metade dos exemplos positivos e, entre eles, reforço do DENTEX e das lesões de Hanói com diâmetro equivalente menor que 35 px. O modelo salvo foi o de melhor média entre a AUC geral, a AUC no DENTEX e a AUC nas lesões pequenas, em 200 imagens de Hanói e 60 do DENTEX reservadas para validação.

Como o classificador dá um ponto por dente, e não um contorno, ele foi comparado com a U-Net v2 por um critério por ponto aplicado igualmente aos dois: uma lesão marcada foi considerada encontrada se algum ponto previsto estava a até 30 px dela, e um ponto previsto foi considerado correto se estava a até 30 px de alguma lesão. Pontos a menos de 25 px entre si foram fundidos. Para a U-Net v2, os pontos foram os picos do mapa de probabilidade suavizado. Também avaliamos a combinação dos dois: os pontos da U-Net v2 no seu limiar de operação somados aos do classificador.

\subsection{Detecção e métricas por lesão}

O mapa de probabilidade de cada modelo foi convertido em lesões pelo mesmo procedimento. O mapa foi suavizado (gaussiana, $\sigma = 4$ px), e cada máximo local em uma janela de $31 \times 31$ pixels com probabilidade $\geq 0{,}05$ originou uma candidata. A região da candidata foi o conjunto conectado de pixels acima de metade do pico, limitado a 40 pixels do pico para o Random Forest e a 100 pixels para as redes, cujo mapa é mais nítido. A confiança de cada candidata foi o valor do pico.

Uma região prevista foi considerada correta quando ao menos 20\% de sua área caía sobre uma lesão marcada; uma lesão marcada foi considerada encontrada quando tocada por uma região correta. Variando o limiar de confiança de 0,05 a 0,95, traçamos curvas FROC, isto é, a sensibilidade por lesão em função dos falsos positivos por imagem [14]. O limiar de operação de cada modelo foi o de maior F1 entre sensibilidade e precisão. Como esse limiar foi escolhido nas mesmas 20 imagens, os valores no ponto de operação são levemente otimistas. Relatamos também o Dice entre os pixels previstos e marcados, somados sobre as imagens. Intervalos de confiança de 95\% foram calculados pelo método de Wilson para proporções [15] e por \textit{bootstrap} por paciente (2.000 reamostragens) para o Dice.

\subsection{Protocolo de validação}

O RF B, as U-Nets v1 e v2 e as duas etapas não usaram nenhuma das 20 imagens locais no treinamento; para eles, as 20 imagens são inteiramente de teste. O RF C e a U-Net com ajuste fino usaram as imagens locais e foram avaliados por validação cruzada em cinco partes fixas de quatro pacientes: cada paciente foi previsto por um modelo treinado sem ele.

\subsection{Aplicação e ferramenta de revisão}

O modelo escolhido foi aplicado às 422 panorâmicas sem marcação. Para cada imagem, gravamos a máscara na resolução original, a radiografia com as lesões contornadas e uma tabela com posição, tamanho e confiança de cada lesão. Para medir a precisão nesse conjunto, desenvolvemos uma ferramenta local de revisão. Nela, o especialista confirma, rejeita ou marca como duvidosa cada detecção, desenha as lesões não detectadas e registra imagens inadequadas; as detecções podem ser ocultadas para que a leitura não seja induzida pelo modelo. As revisões geram estimativas de precisão e sensibilidade com IC95\% e novos rótulos de treinamento; nenhuma revisão havia sido concluída na redação deste rascunho.

\subsection{Software e privacidade}

O pré-processamento, o Random Forest e a avaliação foram implementados em Python 3.11 com NumPy, SciPy, OpenCV, scikit-image e scikit-learn 1.5 [16]; as redes, em PyTorch [17] e torchvision. As figuras quantitativas foram produzidas em R 4.6 com ggplot2 e patchwork. Somente os conjuntos públicos de Hanói e DENTEX foram enviados ao Google Colab (GPU NVIDIA T4). As radiografias de pacientes foram processadas, avaliadas e revisadas exclusivamente no computador local, em CPU de quatro núcleos: cerca de 7,6 s por panorâmica com a U-Net v1, 13--14 s com a U-Net v2 e 19 s com as duas etapas.

\section{Resultados}

\subsection{As redes encontraram quatro a seis vezes mais lesões que o Random Forest}

Comparamos as curvas FROC dos modelos nas 45 lesões (Fig.~3a, Tabela~1). O melhor Random Forest encontrou no máximo 16\% das lesões (RF B; 7/45; IC95\% 8--29\%), com 5,8 falsos positivos por imagem nesse ponto, e nenhuma variante passou de 11\% com até 2 falsos positivos por imagem. A inclusão das imagens locais no treinamento não mudou esse quadro (RF C: 11\% de sensibilidade, 25\% de precisão e 0,6 falso positivo por imagem). Os mapas explicam a falha (Fig.~4): o Random Forest atribuiu probabilidade intermediária a grandes extensões de osso trabecular, e os três atributos de maior importância no RF C foram as posições vertical e horizontal e a distância ao centro, que somaram 31\% da importância por redução de impureza de Gini. O classificador aprendeu onde as lesões costumam estar, e não como elas se apresentam.

A U-Net v1, treinada apenas com Hanói, encontrou 64\% das lesões no ponto de operação (29/45; IC95\% 50--77\%) e 69\% no limiar mais permissivo, sem falsos positivos em nenhuma das 20 imagens (precisão 28/28; IC95\% 88--100\%), com Dice de 0,57. O ajuste fino com as imagens locais acrescentou uma lesão (30/45; 67\%) e elevou o Dice a 0,60, sem mudar a sensibilidade máxima (69\%).

\widefigure{145mm}{fig1_validacao.pdf}{\textbf{Validação e escolha do modelo nas 20 panorâmicas locais (45 lesões).} \textbf{a}, curvas FROC (maior sensibilidade alcançada com cada número de falsos positivos por imagem); pontos, limiar de operação (melhor F1). O RF B, cuja curva fica abaixo da do RF C até 4 falsos positivos por imagem, foi omitido. \textbf{b}, sensibilidade, precisão e Dice no limiar de operação; barras de erro, IC95\% (Wilson para proporções; \textit{bootstrap} por paciente para o Dice). Cores como em \textbf{a}. \textbf{c}, cada coluna é uma lesão marcada, ordenada por metade da imagem e diâmetro equivalente (faixa cinza inferior, do menor ao maior); células coloridas, lesão encontrada pelo modelo no limiar de operação. \textbf{d}, treinamento no Colab: Dice na validação interna de Hanói da U-Net v1, F1 dos ápices (até 20 px) na validação do DENTEX para o localizador e Dice nos recortes de validação para o segmentador da segunda etapa.}

\begin{table*}[!t]
\caption{\textbf{Desempenho por lesão nas 20 panorâmicas locais (45 lesões).} Sensibilidade com no máximo 0,5, 1, 2 e 4 falsos positivos (FP) por imagem e, à direita, valores no limiar de operação (melhor F1), com IC95\% de Wilson ou \textit{bootstrap}. RF C e as U-Nets com ajuste fino foram avaliados por validação cruzada em cinco partes; os demais não usaram as imagens locais.}
\centering\tablefont
\begin{tabularx}{\textwidth}{L *{4}{c} Y Y c Y}
\toprule
Modelo & 0,5 FP & 1 FP & 2 FP & 4 FP & Sensibilidade & Precisão & FP/imagem & Dice \\
\midrule
RF B (Hanói) & 4\% & 4\% & 7\% & 11\% & 16\% (8--29) & 6\% (3--11) & 5,8 & 0,04 (0,02--0,06) \\
RF C (Hanói + locais) & 4\% & 11\% & 11\% & 11\% & 11\% (5--23) & 25\% (10--49) & 0,6 & 0,14 (0,06--0,25) \\
U-Net v1 (Hanói, 1.400 px) & 69\% & 69\% & 69\% & 69\% & 64\% (50--77) & \textbf{100\% (88--100)} & \textbf{0,0} & 0,57 (0,52--0,62) \\
U-Net v1 + ajuste fino & 69\% & 69\% & 69\% & 69\% & 67\% (52--79) & \textbf{100\% (89--100)} & \textbf{0,0} & 0,60 (0,55--0,67) \\
U-Net v2 (2.000 px, Tversky) & \textbf{76\%} & \textbf{76\%} & \textbf{76\%} & \textbf{76\%} & \textbf{69\% (54--80)} & \textbf{100\% (88--100)} & \textbf{0,0} & 0,58 (0,55--0,61) \\
U-Net v2 + ajuste fino & 71\% & \textbf{76\%} & \textbf{76\%} & \textbf{76\%} & 67\% (52--79) & 97\% (83--99) & 0,05 & \textbf{0,62 (0,58--0,66)} \\
Duas etapas (ápices + recortes) & 64\% & 67\% & 69\% & 69\% & 64\% (50--77) & 76\% (61--87) & 0,45 & 0,53 (0,48--0,58) \\
\bottomrule
\end{tabularx}
\end{table*}

\subsection{A U-Net v2 foi o melhor modelo}

A U-Net v2, treinada a 2.000 px com perda de Tversky, encontrou 69\% das lesões no ponto de operação (31/45; IC95\% 54--80\%), sem nenhum falso positivo nas 20 imagens (precisão 29/29), com F1 de 0,82 e Dice de 0,58 (Tabela~1). Aceitando até 0,5 falso positivo por imagem, encontrou 76\% (34/45), a maior sensibilidade entre todos os modelos. O ganho sobre a v1 veio sobretudo da maxila: a v2 encontrou 12 das 20 lesões maxilares e 19 das 25 mandibulares (Fig.~3c). Em 10 dos 20 pacientes, todas as lesões marcadas foram encontradas (Fig.~5). Os intervalos de confiança da v1 e da v2 se sobrepõem, e a diferença no ponto de operação é de duas lesões; a v2 foi escolhida para a aplicação por ter a maior sensibilidade em todos os níveis de falso positivo sem perder precisão e por não depender das imagens locais. O ajuste fino da v2 com as imagens locais produziu o melhor contorno de todos os modelos (Dice 0,62), mas não encontrou mais lesões: 67\% no ponto de operação e os mesmos 76\% a partir de 1 falso positivo por imagem.

\begin{figure*}[!t]
\centering
\includegraphics[width=\textwidth]{fig4_rf_vs_unet.png}
\caption{\textbf{Mapas de probabilidade e detecções em três pacientes.} Primeira coluna: radiografia com as lesões marcadas pelo especialista (verde). Demais colunas: probabilidade de lesão (azul, de transparente em 0 a opaco em 1), lesões marcadas (verde) e detecções no limiar de operação (ciano) para o RF C, a U-Net v2 e o método em duas etapas. Na última coluna, os quadrados brancos são as janelas de recorte centradas nos ápices detectados pelo localizador; fora delas a probabilidade é zero. O RF C foi previsto por modelos que não viram o paciente; as redes não viram nenhuma imagem local. Em \textbf{b}, lesões maxilares pequenas não foram encontradas por nenhum modelo.}
\end{figure*}

\subsection{O método em duas etapas localizou as lesões, mas não as reconheceu melhor}

O localizador de ápices pôs 93\% das lesões marcadas (42/45) dentro de alguma janela de recorte: 19 das 20 maxilares e 23 das 25 mandibulares. Esse é o teto de sensibilidade da segunda etapa, bem acima dos 76\% da U-Net v2. Ainda assim, o método em duas etapas encontrou 64\% das lesões com 0,45 falso positivo por imagem e 69\% com 2, com precisão de 76\% e Dice de 0,53 no ponto de operação (Tabela~1). Na maxila, encontrou 9 das 20 lesões, menos que a U-Net v2.

As duas abordagens encontraram em grande parte as mesmas lesões (Fig.~3c). Das 45, 28 foram encontradas por ambas, 3 apenas pela U-Net v2 e 1 apenas pelas duas etapas; 13 não foram encontradas por nenhuma delas. Nenhum dos modelos avaliados, incluindo o Random Forest, encontrou essas 13 lesões. Elas eram pequenas (diâmetro equivalente mediano de 19 px, contra 41 px das encontradas pela v2) e majoritariamente maxilares (8 das 13). Como 11 delas estavam dentro de uma janela de recorte, examinadas a duas vezes a resolução de trabalho, a falha não decorre de a lesão estar fora da área examinada nem de pouca resolução: o segmentador, treinado com as lesões de Hanói, não as reconhece.

\subsection{O classificador por dente também não reconheceu as lesões difíceis}

Nos dados públicos, o classificador por dente separou muito bem ápices com e sem lesão: AUC de 0,98 nas imagens reservadas de Hanói e do DENTEX, 0,998 só no DENTEX e 0,96 nas lesões pequenas de Hanói (passo 4.000). Nas 20 panorâmicas locais, porém, encontrou no máximo 67\% das lesões, com 4 falsos positivos por imagem, e seu melhor F1 foi 0,55 (Tabela~2). Pelo mesmo critério por ponto, a U-Net v2 manteve 76\% e F1 de 0,82. Somado à U-Net v2, o classificador acrescentou duas lesões (36/45; 80\%), mas só com cerca de 1 a 4 falsos positivos por imagem; uma delas fica a 8 px de outra lesão que a U-Net v2 já encontrava, de modo que parte do ganho pode vir do critério de distância. Na maxila, o classificador encontrou 9 das 20 lesões, contra 12 da U-Net v2. A queda entre a validação pública e as imagens locais indica que as lesões difíceis deste serviço diferem das lesões dos conjuntos públicos.

\begin{table}[!t]
\caption{\textbf{Classificador por dente e U-Net v2 pelo critério por ponto (45 lesões).} Lesão encontrada: ponto previsto a até 30 px dela. Sensibilidade com no máximo 1, 2 e 4 falsos positivos (FP) por imagem e, à direita, o ponto de melhor F1.}
\centering\tablefont
\begin{tabularx}{\columnwidth}{L c c c c}
\toprule
Modelo & 1 FP & 2 FP & 4 FP & Melhor F1 \\
\midrule
U-Net v2 & \textbf{76\%} & \textbf{76\%} & 76\% & \textbf{0,82} (33/45) \\
Classificador por dente & 44\% & 62\% & 67\% & 0,55 (25/45) \\
U-Net v2 + classificador & \textbf{76\%} & \textbf{76\%} & \textbf{80\%} & 0,75 (34/45) \\
\bottomrule
\end{tabularx}
\end{table}

\subsection{Parte das lesões perdidas pode não ser lesão distinta}

A inspeção das 14 lesões que a U-Net v2 não encontrou mostrou dois grupos que exigem revisão do especialista antes de serem tratados como erros do modelo. Seis estavam a até 40 px de outra lesão marcada que a rede encontrou (pacientes 03, 09, 12, 13 e duas no paciente 20) e podem ser parte da mesma lesão, desenhada em dois traços. Quatro tinham diâmetro equivalente de 12 px ou menos, menor que um ápice radicular (duas no paciente 02, uma no 15 e uma no 20), e podem ser resíduos do traço. Somados, sem repetição, são 9 casos. Se todos fossem excluídos da referência, a sensibilidade da U-Net v2 no ponto de operação seria de 31/36 (86\%); trata-se de um limite superior, e não de um resultado, até que o especialista julgue cada caso. Para isso, preparamos uma página local que mostra cada lesão perdida ampliada, com o traço original, a máscara derivada e o mapa de probabilidade.

\begin{figure*}[p]
\centering
\includegraphics[width=\textwidth,height=215mm,keepaspectratio]{fig5_validacao_20_pacientes.png}
\caption{\textbf{Detecções da U-Net v2 nas 20 panorâmicas locais.} Verde: lesões marcadas pelo especialista; ciano: detecções no limiar de operação (0,5). O título de cada painel indica o paciente, as lesões encontradas e os falsos positivos, quando houver. Nenhuma das imagens foi usada no treinamento.}
\end{figure*}

\subsection{Limiares e escalas não recuperaram as lesões perdidas}

Antes de treinar a v2, testamos se as lesões perdidas pela U-Net v1 estavam logo abaixo do limiar ou fora da escala de treinamento (Tabela~3). Mesmo aceitando qualquer pico acima de 0,002, a sensibilidade chegou a 73\%, com 0,75 falso positivo por imagem. Prever também na imagem ampliada 1,5 vez encontrou duas lesões adicionais ao custo de cerca de um falso positivo por imagem, e a ampliação de 2 vezes levou a sensibilidade a 7\%. Ampliar a imagem apenas na inferência não substitui treinar na escala maior, o que motivou a v2; a v2 e as duas etapas, porém, também não recuperaram as 13 lesões.

\begin{table*}[!t]
\caption{\textbf{Análises de sensibilidade da U-Net v1 nas 20 panorâmicas locais.} Os limiares e as combinações foram examinados nas mesmas imagens; os resultados descrevem o comportamento do modelo e não foram usados para escolher a versão aplicada.}
\centering\tablefont
\begin{tabularx}{\textwidth}{>{\RaggedRight\arraybackslash}p{52mm} c c L}
\toprule
Variante & Lesões encontradas & FP/imagem & Leitura \\
\midrule
Limiar de operação (0,5) & 29/45 (64\%) & 0,00 & referência \\
Limiar 0,2 & 31/45 (69\%) & 0,20 & duas lesões a mais \\
Limiar e pico mínimo 0,002 & 33/45 (73\%) & 0,75 & ganho pequeno, com quase um falso positivo por imagem \\
Imagem e ampliação 1,5$\times$ (máximo) & 33/45 (73\%) & $\leq$1,0 & F1 0,76, contra 0,78 na escala original \\
Somente ampliação 1,5$\times$ & 28/45 (62\%) & $\leq$1,0 & piora \\
Somente ampliação 2$\times$ & 3/45 (7\%) & $\leq$0,5 & a rede não reconhece lesões fora da escala de treino \\
\bottomrule
\end{tabularx}
\end{table*}

\subsection{Aplicação a 422 panorâmicas sem marcação}

A U-Net v2 assinalou 807 lesões em 407 das 422 panorâmicas (96\%); 15 imagens não tiveram detecção, e 47 tiveram quatro ou mais (Fig.~6a, Fig.~7). A U-Net v1 havia assinalado 673 lesões em 376 imagens; o aumento de 20\% no número de detecções acompanha a maior sensibilidade da v2, mas também pode incluir mais falsos positivos. As detecções concentraram-se junto aos ápices das duas arcadas, com 64\% delas na metade inferior da imagem (Fig.~6b). O diâmetro equivalente mediano das lesões detectadas foi 39 pixels (intervalo interquartil 32--54), próximo ao das lesões encontradas na validação e maior que o das marcadas em geral (32 pixels), o que sugere que lesões pequenas também foram perdidas nesse conjunto (Fig.~6c). A confiança foi $\geq 0{,}95$ em 89\% das detecções (Fig.~6d). Sem marcação de referência, a precisão nesse conjunto não pode ser medida, e a proporção de imagens com detecção ainda precisa ser confirmada pela revisão especializada.

\widefigure{100mm}{fig2_aplicacao.pdf}{\textbf{Aplicação da U-Net v2 às 422 panorâmicas sem marcação.} \textbf{a}, número de lesões detectadas por imagem. \textbf{b}, posição dos centros das detecções (fração da largura e da altura da imagem); manchas, densidade bidimensional por \textit{kernel}; as linhas tracejadas são referências aproximadas, não anatômicas. \textbf{c}, distribuição acumulada do diâmetro equivalente das lesões detectadas na aplicação e das lesões marcadas na validação, encontradas e não encontradas; tracejados, medianas. \textbf{d}, confiança das detecções (pico do mapa de probabilidade suavizado).}

\begin{figure*}[!t]
\centering
\includegraphics[width=\textwidth]{fig7_aplicacao_exemplos.png}
\caption{\textbf{Exemplos de detecções da U-Net v2 no acervo sem marcação.} Uma imagem sorteada em cada faixa de número de detecções (0, 1, 2--3 e 4 ou mais). Os identificadores (A001--A422) são anônimos. As detecções não foram confirmadas por especialista.}
\end{figure*}

\scopenote{\textbf{Escopo da evidência.} As métricas de validação vêm de 45 lesões em 20 panorâmicas de um único serviço, marcadas por \pendente{número de} examinador(es), com limiar de operação escolhido nas mesmas imagens. Não há leitura repetida para medir a variabilidade entre examinadores. As diferenças entre as redes (uma a cinco lesões) estão dentro dos intervalos de confiança. Os resultados no acervo de 422 imagens descrevem o que a U-Net v2 assinalou, e não sua precisão. Nenhum resultado deste rascunho é diagnóstico.}

\section{Discussão}

O resultado central é que redes convolucionais treinadas apenas com panorâmicas públicas de outro país, rotuladas por caixas, encontraram dois terços a três quartos das lesões periapicais marcadas em um serviço local, praticamente sem falsos positivos, enquanto um classificador por pixel treinado com os mesmos dados não passou de 16\%. A diferença decorre do que cada modelo consegue representar: a lesão periapical é uma radiolucidez definida pela relação com uma raiz específica, e não uma textura local. O Random Forest, mesmo com atributos de contexto e negativos difíceis, respondeu à textura trabecular e à posição; a U-Net integra o contexto em várias escalas. Essa conclusão contrasta com o sucesso de classificadores por pixel em problemas nos quais o alvo tem assinatura local própria [5,6].

O segundo resultado é que quatro intervenções plausíveis para as lesões pequenas tiveram efeito limitado. Treinar em resolução maior com uma perda que pesa a lesão perdida acrescentou duas lesões no ponto de operação e três na sensibilidade máxima, com o ganho concentrado na maxila (10 para 12 de 20). O ajuste fino local acrescentou uma. O método em duas etapas, que examinou a vizinhança de cada ápice a duas vezes a resolução de trabalho e cobriu 93\% das lesões, não superou a U-Net v2. O classificador por dente, com o dobro de contexto, reforço de lesões pequenas e rótulos de outro aparelho, foi excelente nos dados públicos (AUC 0,98) e fraco nos locais. Das 13 lesões que nenhum modelo encontrou, 11 estavam dentro da área examinada pela segunda etapa. O gargalo, portanto, não é localizar a região nem enxergá-la com resolução suficiente, e sim reconhecer a lesão pequena e maxilar, um padrão pouco representado nas caixas de Hanói ou diferente delas em aparência. Na maxila, a sobreposição do seio maxilar e do palato reduz ainda mais o contraste.

Essa leitura tem uma consequência prática: o próximo ganho deve vir de dados, e não de arquitetura. Antes disso, porém, a própria referência precisa ser revista: nove das 14 lesões perdidas são marcas pequenas ou vizinhas de lesões encontradas, e a sensibilidade real pode estar entre 69\% e 86\%. Exemplos locais de lesões pequenas e maxilares, obtidos com a ferramenta de revisão, são o meio mais direto de ensinar à rede o padrão que falta. Vinte imagens com 45 lesões não bastaram no ajuste fino, mas um conjunto maior, de algumas centenas de lesões revisadas, permitiria testar essa hipótese. Rótulos mais fiéis também ajudariam: as elipses inscritas nas caixas fazem as redes desenharem regiões arredondadas, o que limita o Dice. Modelos de segmentação guiados por caixa, como o MedSAM [19], poderiam converter as caixas de Hanói em contornos sem nova marcação manual.

O método em duas etapas ainda tem utilidade. Seu localizador, treinado sem nenhuma lesão, transforma a panorâmica em cerca de 30 regiões apicais padronizadas, sobre as quais um especialista poderia revisar rapidamente a probabilidade de lesão dente a dente. Sua menor precisão (76\%, contra 100\% da U-Net v2) vem de recortes que veem pouco contexto e, por isso, confundem estruturas radiolúcidas normais com lesões. Combinar o mapa da U-Net v2, que tem contexto, com o da segunda etapa, que tem resolução, encontraria no máximo 32 das 45 lesões nesta validação, uma a mais que a v2.

A sensibilidade de 69--76\% deve ser lida com cuidado. Ela foi medida em 45 lesões de 20 pacientes, com intervalo de confiança de cerca de 25 pontos percentuais, e o limiar foi escolhido nas mesmas imagens. A referência é a marcação de \pendente{número de} examinador(es) sobre panorâmicas, exame cuja sensibilidade é limitada em relação à tomografia [1]. A comparação com a taxa média de verdadeiros positivos de 0,51 relatada para cirurgiões em panorâmicas [2] é sugestiva, mas não direta, porque populações, referências e critérios de acerto diferem.

No acervo de 422 imagens, a ausência de referência impede afirmar acurácia. A proporção de imagens com ao menos uma detecção pode refletir a população atendida, falsos positivos não observados na validação ou ambos. A revisão de uma amostra aleatória dessas imagens, e não apenas das prioritárias, transforma a aplicação em evidência de desempenho; duas leituras independentes de um subconjunto comum medirão a concordância entre examinadores, teto realista para qualquer modelo. Também há limitações de dados: a separação entre maxila e mandíbula foi aproximada pela metade da imagem, o acervo contém possíveis duplicatas de pacientes, uma das 45 máscaras incluiu a raiz porque o preenchimento recorreu ao fecho convexo, e os ápices do DENTEX foram derivados por regra geométrica, não marcados.

\section{Conclusão}

Com apenas 20 panorâmicas marcadas no próprio serviço, uma U-Net treinada em panorâmicas públicas a 2.000 px com perda de Tversky encontrou 69\% das lesões periapicais sem falsos positivos e 76\% com até 0,5 falso positivo por imagem, enquanto um Random Forest por pixel não passou de 16\%. Um método em duas etapas guiado por ápices cobriu 93\% das lesões com seus recortes, e um classificador por dente atingiu AUC de 0,98 nos dados públicos, mas nenhum dos dois reconheceu melhor as lesões locais. As lesões perdidas por todos os modelos foram pequenas e sobretudo maxilares, o que aponta para dados de treinamento, e não para arquitetura ou resolução, como o limite atual. A evidência justifica usar a U-Net v2 como apoio à triagem de lesões evidentes, sob revisão especializada. Ela não justifica afirmar desempenho diagnóstico, precisão no acervo sem marcação ou generalização para outros serviços.

\section*{Aprovação ética e consentimento}

\pendente{Informar o comitê de ética em pesquisa, o número do parecer (CAAE) e se houve consentimento dos pacientes ou dispensa formal, inclusive para a publicação das imagens.} As radiografias de pacientes foram processadas somente em computador local; nenhuma imagem de paciente foi enviada a serviços de nuvem. As figuras identificam pacientes apenas por números sequenciais.

\section*{Disponibilidade de dados e código}

Os conjuntos de Hanói [7] e DENTEX [8] estão disponíveis publicamente; o DENTEX sob licença não comercial. As radiografias locais e as marcações não são públicas. \pendente{Informar a forma de acesso e as restrições.} O código de pré-processamento, treinamento, avaliação, figuras e revisão está no repositório do projeto \pendente{endereço do repositório}.

\section*{Contribuições dos autores}

\pendente{Descrever as contribuições segundo a taxonomia CRediT.}

\section*{Declaração de conflitos de interesse}

\pendente{Declarar conflitos de interesse ou a ausência deles.}

\section*{Uso de inteligência artificial}

Uma ferramenta de inteligência artificial generativa foi usada para implementar e executar as análises e para redigir este rascunho, a partir de dados, decisões e instruções do autor. Todas as afirmações, números, referências e interpretações devem ser verificados e aprovados pelo autor, que permanece responsável pelo conteúdo final.

\balance
\section*{Referências}

{\sffamily\fontsize{7.8pt}{10.2pt}\selectfont
\begin{enumerate}[leftmargin=5mm,itemsep=1mm]
\item Estrela, C., Bueno, M. R., Leles, C. R., Azevedo, B. \& Azevedo, J. R. Accuracy of cone beam computed tomography and panoramic and periapical radiography for detection of apical periodontitis. \textit{Journal of Endodontics} \textbf{34}, 273-279 (2008). \href{https://doi.org/10.1016/j.joen.2007.11.023}{doi:10.1016/j.joen.2007.11.023}
\item Endres, M. G. \textit{et al.} Development of a deep learning algorithm for periapical disease detection in dental radiographs. \textit{Diagnostics} \textbf{10}, 430 (2020). \href{https://doi.org/10.3390/diagnostics10060430}{doi:10.3390/diagnostics10060430}
\item Ekert, T. \textit{et al.} Deep learning for the radiographic detection of apical lesions. \textit{Journal of Endodontics} \textbf{45}, 917-922.e5 (2019).
\item Breiman, L. Random forests. \textit{Machine Learning} \textbf{45}, 5-32 (2001). \href{https://doi.org/10.1023/A:1010933404324}{doi:10.1023/A:1010933404324}
\item Berg, S. \textit{et al.} ilastik: interactive machine learning for (bio)image analysis. \textit{Nature Methods} \textbf{16}, 1226-1232 (2019). \href{https://doi.org/10.1038/s41592-019-0582-9}{doi:10.1038/s41592-019-0582-9}
\item Arganda-Carreras, I. \textit{et al.} Trainable Weka Segmentation: a machine learning tool for microscopy pixel classification. \textit{Bioinformatics} \textbf{33}, 2424-2426 (2017). \href{https://doi.org/10.1093/bioinformatics/btx180}{doi:10.1093/bioinformatics/btx180}
\item Do, H. V. \textit{et al.} A dataset of apical periodontitis lesions in panoramic radiographs for deep-learning-based classification and detection. \textit{Data in Brief} \textbf{54}, 110486 (2024). Dados: Panoramic radiographs with periapical lesions, Mendeley Data v3, \href{https://doi.org/10.17632/kx52tk2ddj.3}{doi:10.17632/kx52tk2ddj.3}
\item Hamamci, I. E. \textit{et al.} DENTEX: an abnormal tooth detection with dental enumeration and diagnosis benchmark for panoramic X-rays. arXiv:2305.19112 (2023).
\item Ronneberger, O., Fischer, P. \& Brox, T. U-Net: convolutional networks for biomedical image segmentation. In \textit{Medical Image Computing and Computer-Assisted Intervention -- MICCAI 2015}, LNCS 9351, 234-241 (2015). \href{https://doi.org/10.1007/978-3-319-24574-4_28}{doi:10.1007/978-3-319-24574-4\_28}
\item Ørstavik, D., Kerekes, K. \& Eriksen, H. M. The periapical index: a scoring system for radiographic assessment of apical periodontitis. \textit{Endodontics \& Dental Traumatology} \textbf{2}, 20-34 (1986). \href{https://doi.org/10.1111/j.1600-9657.1986.tb00119.x}{doi:10.1111/j.1600-9657.1986.tb00119.x}
\item Zuiderveld, K. Contrast limited adaptive histogram equalization. In \textit{Graphics Gems IV} (ed. Heckbert, P. S.) 474-485 (Academic Press, 1994).
\item He, K., Zhang, X., Ren, S. \& Sun, J. Deep residual learning for image recognition. In \textit{Proceedings of the IEEE Conference on Computer Vision and Pattern Recognition}, 770-778 (2016). \href{https://doi.org/10.1109/CVPR.2016.90}{doi:10.1109/CVPR.2016.90}
\item Loshchilov, I. \& Hutter, F. Decoupled weight decay regularization. In \textit{International Conference on Learning Representations} (2019).
\item Chakraborty, D. P. Maximum likelihood analysis of free-response receiver operating characteristic (FROC) data. \textit{Medical Physics} \textbf{16}, 561-568 (1989). \href{https://doi.org/10.1118/1.596358}{doi:10.1118/1.596358}
\item Wilson, E. B. Probable inference, the law of succession, and statistical inference. \textit{Journal of the American Statistical Association} \textbf{22}, 209-212 (1927). \href{https://doi.org/10.1080/01621459.1927.10502953}{doi:10.1080/01621459.1927.10502953}
\item Pedregosa, F. \textit{et al.} Scikit-learn: machine learning in Python. \textit{Journal of Machine Learning Research} \textbf{12}, 2825-2830 (2011).
\item Paszke, A. \textit{et al.} PyTorch: an imperative style, high-performance deep learning library. In \textit{Advances in Neural Information Processing Systems 32} (2019).
\item Salehi, S. S. M., Erdogmus, D. \& Gholipour, A. Tversky loss function for image segmentation using 3D fully convolutional deep networks. In \textit{Machine Learning in Medical Imaging}, LNCS 10541, 379-387 (2017). \href{https://doi.org/10.1007/978-3-319-67389-3_44}{doi:10.1007/978-3-319-67389-3\_44}
\item Ma, J. \textit{et al.} Segment anything in medical images. \textit{Nature Communications} \textbf{15}, 654 (2024). \href{https://doi.org/10.1038/s41467-024-44824-z}{doi:10.1038/s41467-024-44824-z}
\end{enumerate}}

\end{document}
